\documentclass[11pt]{article}
\usepackage{amsmath, amssymb, amsthm}
\usepackage{geometry}
\usepackage{hyperref}
\usepackage{graphicx}
\usepackage{bm}
\usepackage{cite}

\geometry{margin=1in}

\title{\textbf{The Gravitational Baseline Principle: Structural Asymmetry and the Pre-Geometric Quantum Substrate}}
\author{Paul Jarvis \\ Independent Researcher, UK \\ \texttt{mrpaulwjarvis@gmail.com}}
\date{June 6, 2026}

\begin{document}

\maketitle

\begin{abstract}
We formalize the conceptual incompatibility between General Relativity (GR) and Quantum Field Theory (QFT) by identifying a fundamental structural asymmetry: standard quantum interactions represent local operator fields acting within a kinematic background, whereas gravitational dynamics dictate the background metric structure itself. We state this architectural divergence as the \emph{Gravitational Baseline Principle}. We demonstrate that treating gravity as a conventional, quantizable gauge field introduces unavoidable mathematical contradictions regarding background independence and gauge covariance. To resolve this, we propose a minimal, background-independent mathematical framework in which both smooth Riemannian manifolds and quantum field algebras co-emerge from a discrete, relational quantum substrate governed by density-matrix dynamics. In this framework, the graviton is recovered not as an elementary gauge boson, but as a collective, low-energy excitation of the pre-geometric medium.
\end{abstract}

\section{Introduction}

The century-long impasse in formulating a consistent theory of quantum gravity is not a crisis of empirical data; it is fundamentally a crisis of structural architecture. Quantum Field Theory (QFT) describes matter and radiation via operator-valued distributions evolving on a fixed, smooth kinematic background. Conversely, General Relativity (GR) denies the existence of any prior background container, asserting that spacetime geometry is a dynamical physical degree of freedom governed by the distribution of stress-energy.

In this paper, we argue that the failure of standard covariant quantization techniques stems from an unaddressed asymmetry. Gravity is not a selective force acting upon a subset of fields; it is the universal geometric framework that defines the causal and metrological relationships between all fields. We formalize this insight via the \emph{Gravitational Baseline Principle} and outline the minimal mathematical constraints required for a pre-geometric quantum substrate to successfully generate smooth continuous physical manifolds.

\section{The Structural Asymmetry Problem}

Einstein's field equations,
\begin{equation}
G_{\mu\nu} = 8\pi G\, T_{\mu\nu},
\end{equation}
dictate that all physical fields couple universally to the metric tensor $g_{\mu\nu}$. In contrast, the gauge interactions of the Standard Model operate exclusively on specific representations of internal symmetry groups.

This establishes a profound architectural conflict when standard perturbative quantization techniques are applied to gravity. QFT constructs a Fock space by defining creation and annihilation operators relative to a fixed global vacuum state, which requires a background metric to define microcausality. If the metric is itself perturbed and quantized as a field $h_{\mu\nu}$, the framework violates diffeomorphism invariance at the foundational level.

\section{The Gravitational Baseline Principle}

Gravity is not a peer field to the gauge forces within a spacetime container, but the geometric framework that renders gauge field relationships definable. Because all physical systems are non-linearly coupled to this background, gravity serves as the absolute measurement baseline. Unification cannot occur by expanding the operator algebra of fields on a background, but must proceed from a pre-geometric layer where spacetime geometry is non-fundamental.

The Planck scales,
\begin{equation}
\ell_P = \sqrt{\frac{\hbar G}{c^3}}, \qquad t_P = \sqrt{\frac{\hbar G}{c^5}},
\end{equation}
signify the physical scale at which the continuum description of geometry breaks down entirely, exposing the pre-geometric substrate.

\section{Formal Substrate Architecture}

We define the pre-geometric substrate $Q$ as an abstract quantum system described by a state vector $|\Psi\rangle$ in a non-local, background-independent Hilbert space $H_Q$. The geometric connectivity of the emergent space is encoded directly within the entanglement structure of the state.

Let $\rho_0 = |\Psi\rangle\langle\Psi|$ be the density matrix of the substrate. We extract a relational network profile by evaluating the mutual information:
\begin{equation}
I(i,j) = S(\rho_i) + S(\rho_j) - S(\rho_{ij}),
\end{equation}
yielding a symmetric relational connectivity matrix $A_{ij}$.

\subsection{Emergent Continuum Limit}

(Your original derivation preserved here.)

\subsection{Dynamical Law of the Pre-Geometric Substrate}

To complete the definition of the substrate, we introduce a minimal, background-independent dynamical law for $\rho_0$:
\begin{equation}
\frac{d\rho_0}{dt} = -i[H_Q, \rho_0].
\end{equation}
The Hamiltonian is constructed from the relational matrix:
\begin{equation}
H_Q = \sum_{i,j} K(A_{ij})\, \sigma_i \otimes \sigma_j,
\end{equation}
where $K$ is a monotonic kernel of entanglement strength. This ensures non-local microscopic dynamics, locality emerging only when $A_{ij}$ becomes band-limited, and metric structure arising from relaxation of $\rho_0$.

\section{Theoretical Constraints and Peer Review Targets}
\label{sec:constraints}

In this section we outline how the Gravitational Baseline Principle constrains the low-energy limit of the theory, the status of the graviton, and the expected phenomenology. The emphasis is on structural consistency with standard tensor calculus and known properties of linearized gravity.

\subsection{The Status of the Graviton}
\label{subsec:graviton}

Within the present framework, the graviton is not introduced as a fundamental gauge boson, but as an emergent, collective excitation of the underlying relational substrate. At the level of the emergent continuum, we consider small perturbations of the metric around a background solution $g_{\mu\nu}^{(0)}$,
\begin{equation}
g_{\mu\nu}(x) = g^{(0)}_{\mu\nu}(x) + h_{\mu\nu}(x),
\end{equation}
where $|h_{\mu\nu}| \ll 1$.

The key requirement is that, in the weak-field, low-energy limit, the effective action for $h_{\mu\nu}$ reduces to the Pauli--Fierz form on the chosen background. On a flat background $g^{(0)}_{\mu\nu} = \eta_{\mu\nu}$, the standard quadratic action for a free, massless spin-2 field reads
\begin{align}
S_{\text{PF}}[h] = \int d^4x \, \Big[
& - \tfrac{1}{2} \partial_\lambda h_{\mu\nu} \, \partial^\lambda h^{\mu\nu}
+ \partial_\mu h^{\mu\nu} \, \partial^\lambda h_{\lambda\nu}
- \partial_\mu h^{\mu\nu} \, \partial_\nu h
+ \tfrac{1}{2} \partial_\lambda h \, \partial^\lambda h
\Big],
\label{eq:PauliFierz}
\end{align}
where $h \equiv h^{\mu}{}_{\mu}$ is the trace. This action is invariant under the linearized gauge transformation
\begin{equation}
h_{\mu\nu} \rightarrow h_{\mu\nu} + \partial_\mu \xi_\nu + \partial_\nu \xi_\mu,
\end{equation}
which is the infinitesimal remnant of diffeomorphism invariance.

In the present model, the perturbations $h_{\mu\nu}$ are not fundamental fields but effective variables induced by small variations of the relational matrix $A_{ij}$. Schematically, one may write
\begin{equation}
h_{\mu\nu}(x) \sim \sum_{i,j} \delta A_{ij} \, f^{(i)}_{\mu}(x) f^{(j)}_{\nu}(x),
\label{eq:h-from-A}
\end{equation}
where the $f^{(i)}_{\mu}(x)$ encode the coarse-grained embedding of discrete subsystems into the emergent manifold. The detailed form of these functions is model-dependent, but the requirement is that, after coarse-graining and projection onto the symmetric, transverse, traceless sector, the effective action for $h_{\mu\nu}$ reduces to Eq.~\eqref{eq:PauliFierz} in the appropriate limit.

Thus, the ``graviton'' appears as the quantized normal mode of these collective excitations, with the usual spin-2, massless representation of the Poincaré group in the low-energy, flat-background regime. This is fully compatible with standard perturbative gravity, while preserving the non-fundamental status of the metric at the microscopic level.

\subsection{Thermodynamic Gravitational Dynamics}
\label{subsec:thermo}

If smooth geometry is an aggregate manifestation of a relational matrix web, then Einstein's field equations should emerge as a macroscopic thermodynamic equation of state, in the spirit of Jacobson's derivation of the Einstein equations from local Rindler horizon thermodynamics~\cite{jacobson}. In the present framework, the coarse-grained metric $g_{\mu\nu}$ and its associated curvature tensors are effective variables encoding the large-scale organization of quantum mutual information.

Small-scale fluctuations of the pre-geometric substrate are expected to manifest as higher-order geometric corrections to the Einstein--Hilbert action. At the level of an effective field theory, one may anticipate terms of the schematic form
\begin{equation}
S_{\text{eff}} = \frac{1}{16\pi G} \int d^4x \, \sqrt{-g} \left( R + \alpha R^2 + \beta R_{\mu\nu} R^{\mu\nu} + \cdots \right),
\label{eq:Seff-curvature}
\end{equation}
where $\alpha$ and $\beta$ are dimensionful coefficients determined by the entanglement spectrum and the microscopic scale of the substrate. These corrections are compatible with standard effective field theory expectations for gravity at high curvature, but here they acquire a direct interpretation in terms of fluctuations of the relational connectivity.

Moreover, because geometry is encoded in mutual information, there is a natural upper bound on the information content that can be localized within a given region, reminiscent of holographic entropy bounds. Deviations from the classical Bekenstein--Hawking area law may therefore arise as subleading corrections controlled by the structure of $A_{ij}$.

\subsection{Phenomenological Consequences}
\label{subsec:phenomenology}

Although the present framework is primarily structural and conceptual, it suggests several phenomenological avenues:

\begin{itemize}
\item \textbf{Holographic information-capacity limits:}  
If spatial geometry is an emergent encoding of mutual information, then precision interferometry and high-sensitivity experiments may encounter an irreducible noise floor associated with fluctuations of the underlying relational substrate. The detailed spectrum of such noise would depend on the dynamics of $H_Q$, but it would be constrained by holographic-like scaling.

\item \textbf{Higher-order curvature corrections:}  
The terms in Eq.~\eqref{eq:Seff-curvature} can, in principle, leave imprints in strong-gravity regimes, such as black hole quasi-normal modes, early-universe cosmology, or late-time cosmic acceleration. Existing bounds on higher-curvature operators constrain the allowed magnitude of substrate-induced corrections.

\item \textbf{Gravitational-wave propagation:}  
Collective excitations of the substrate corresponding to gravitational waves may exhibit tiny, scale-dependent modifications to their dispersion relation. A generic parametrization would be
\begin{equation}
\omega^2 = k^2 \left[ 1 + \varepsilon(k) \right],
\end{equation}
where $\varepsilon(k)$ is a small function suppressed by the microscopic scale (e.g.\ the Planck scale or the characteristic relational scale). Current observations from LIGO/Virgo/KAGRA are consistent with $\varepsilon(k) \approx 0$ over the probed frequency range, but future detectors could improve these bounds.
\end{itemize}

These phenomenological considerations are not yet precise predictions, but they indicate where the present framework can be confronted with data and how it might be falsified or constrained.

\subsection{Relation to Existing Approaches}
\label{subsec:relation}

It is useful to briefly situate the Gravitational Baseline Principle within the broader landscape of emergent gravity and quantum gravity proposals:

\begin{itemize}
\item \textbf{Tensor-network and entanglement-geometry approaches:}  
These models also relate geometry to entanglement patterns, but typically assume a fixed underlying graph or lattice. In contrast, the present framework treats the relational matrix $A_{ij}$ as dynamical, with geometry emerging from its evolution.

\item \textbf{AdS/CFT and holographic dualities:}  
Holographic constructions relate a bulk gravitational theory to a boundary quantum field theory, with bulk geometry encoded in boundary entanglement. The present approach is not tied to asymptotic AdS structure or a specific boundary theory; it aims instead at a background-independent, pre-geometric substrate.

\item \textbf{Loop Quantum Gravity and spin networks:}  
Loop Quantum Gravity quantizes geometric operators directly, leading to discrete spectra for areas and volumes. Here, by contrast, geometry is not fundamental but arises from quantum information-theoretic relations; discreteness, if present, is a property of the substrate degrees of freedom rather than of geometric operators themselves.

\item \textbf{Causal set theory:}  
Causal set theory posits a discrete, partially ordered set as the fundamental structure, with spacetime emerging as an approximation. The present model shares the idea of a pre-geometric combinatorial structure but uses quantum mutual information rather than causal order as the primary organizing principle.

\item \textbf{Group field theory and related approaches:}  
Group field theories describe pre-geometric quanta whose condensates can give rise to continuum spacetime. The present framework is conceptually similar in treating spacetime as emergent from a many-body quantum system, but it emphasizes density-matrix dynamics and relational information rather than group-theoretic combinatorics.
\end{itemize}

In all cases, the Gravitational Baseline Principle provides a unifying perspective: gravity is the emergent geometric baseline that renders all other field-theoretic relations definable, rather than a peer field to be quantized on a fixed background.


\section{Conclusion}

The Gravitational Baseline Principle provides a mathematically self-consistent conceptual path forward for quantum gravity. By recognizing that gravity is the geometric framework of measurement rather than an active field, we bypass the contradictions of background-dependent quantization. The emergence of a Riemannian spatial metric from a non-local quantum entanglement density matrix demonstrates that a smooth continuum can arise naturally from an abstract relational substrate.

\begin{thebibliography}{9}

\bibitem{dewitt}
B.~S.~DeWitt, ``Quantum Theory of Gravity,'' \emph{Phys. Rev.} \textbf{160}, 1113 (1967).

\bibitem{dirac}
P.~A.~M.~Dirac, \emph{The Principles of Quantum Mechanics}, Oxford University Press (1930).

\bibitem{einstein}
A.~Einstein, ``The Foundation of the General Theory of Relativity,'' \emph{Annalen der Physik} \textbf{49}, 769 (1916).

\bibitem{jacobson}
T.~Jacobson, ``Thermodynamics of Spacetime: The Einstein Equation of State,'' \emph{Phys. Rev. Lett.} \textbf{75}, 1260 (1995).

\bibitem{susskind}
L.~Susskind, ``ER=EPR, GHZ, and the Consistency of Quantum Measurements,'' \emph{Fortschritte der Physik} \textbf{64}, 72 (2016).

\end{thebibliography}

\end{document}
